\begin{abstract}
Weight-only post-training quantization (PTQ) costs some recognition datasets far more accuracy than others.
We study this dataset dependence with half-quadratic quantization (HQQ) on \V{nDatasets} fine-grained datasets, three backbones, three checkpoint regimes and \V{nConditions} combinations of bit width and group size.
Each quantized model is compared with its own unquantized checkpoint through accuracy retention, the ratio of quantized to baseline accuracy.
Bit width and group size set the level of retention.
In the transition between collapse and full retention, five soybean and cotton cultivar datasets keep less accuracy than the other datasets: at 3-bit weights with group size 128 they retain \V{retSc_w3g128} of their accuracy, against \V{retOther_w3g128}.
Dataset difficulty rankings agree across quantization conditions, with Kendall's $W$ of at least \V{wMin}.
Datasets with small baseline logit margins lose more, and matching the margin distributions of the two groups narrows their retention gap by \V{reductionMin}\% to \V{reductionMax}\%.
Baseline accuracy predicts difficulty almost as well as the margin does, yet datasets with equal baseline accuracy can differ in retention by up to \V{pairsMaxDiff}.
Feature similarity between the baseline and quantized models does not track the accuracy loss.
\end{abstract}
